\def\ChapterCode{PSY}
\chapter
[\CO{[PSY]} Psychiatrie]
{Psychiatrie}
\label{chp:PSY}

\ChapterIntro
{%
~~~
}
{%
\begin{description}
  \item[Maintainer:] Sebastian Gabriel
  \item[Autoren:] Diverse
  \item[Reviewer:] Standard-Reviewprozess
  \item[Version:] {\VarDocVersionTypeName} ({\VarDocVersionTypeDescription})
  \item[SHA1:] {\DocGitCommit}
\end{description}

{\TxTeilkapitelAASS}
}

\clearpage

\Dictum[Otto Pötzl, nach \cite{Jolesch2}]{Der, der den
Schlüssel hat, ist der Arzt.}

\nocite{WunnPsych1,OHCSP7}
\label{topic:unterbringungsgesetz}


\section{\HeadingKeyword{Allgemeines}}

\Text{Einleitung}
{%
\DefinitionInPara{Die Psychiatrie beschäftigt sich
  mit der Prävention, Diagnostik und Therapie psychischer Störungen.}
Die Psychiatrie ist ein sehr kompliziertes und
anspruchsvolles Fachgebiet der Medizin. Sie beschäftigt
sich einerseits mit handfesten, körperlich bedingten
Störungen, die Auswirkungen auf das Denken und die geistige
Leistungsfähigkeit haben. Andererseits hat man es in der
Psychiatrie sehr oft mit Störungen zu tun, von deren
Ursache die Wissenschaft kaum eine Ahnung hat. In der
Psychiatrie ist es auch besonders schwierig, zwischen
\enquote*{gesund} und \enquote*{krank} treffsicher zu
unterscheiden.
}


Grundsätzlich muss aber gesagt werden, dass psychiatrische
Erkrankungen\footnote{Sofern der Begriff
\emph{\enquote{Erkrankung}} überhaupt angebracht ist \ldots}
etwas alltägliches und nichts außergewöhnliches sind. Je
nach Studie leiden bis zu 30\,\% der Bevölkerung an
psychiatrischen Erkrankungen, allen voran an Depressionen
und der Alkoholkrankheit.


\Dictum[Redensart]{If there's no problem, don't fix it.}

\paragraph{Der {Krankheitsbegriff} in der Psychiatrie} Man
trifft oft auch auf nicht klar umrissene Krankheitsbilder, und
auch deren jeweiliger Krankheitswert ist oft umstritten,
sehr vom einzelnen Patienten und seiner Situation abhängig:
Während der eine Mensch an seiner Depression leidet,
verarbeitet ein anderer seine Gefühlserfahrungen zu einem Lied und
sichert sich dank der AKM\footnote{\Z{Autoren, Komponisten
und Musikverleger. Die AKM ist die größte Urheberrechts- und
Verwertungsgesellschaft in Österreich.}} eine
Einnahmequelle. (Subjektiv betrachtet erscheint die
Depression manchmal als ein treibender Motor der
Singer-/Songwriter-Szene.)

Einer der bekanntesten psychiatrischen Patienten war der
niederländische Maler \E{Vincent van Gogh}
({\sffamily\textborn}\,1853\,--\,{\sffamily\textdied}\,1890).
Im Alter von 35 Jahren erkrankte er an einer psychiatrischen Erkrankung die
von Wahnvorstellungen, Albträumen, Depressionen und
wiederholten Selbstmordversuchen geprägt war. Es sollte
nicht bei den Versuchen bleiben, doch bis zu seiner
Selbsttötung am 29. Juli 1890 schuf er hunderte Gemälde und
über tausend Zeichnungen. Heute gilt van Gogh als
bedeutender Mitbegründer und Impulsgeber der modernen
Malerei. \cite{wikiVanGogh}

\begin{Synopsis}
  \SynopsisItem{Die psychiatrische Diagnose verhält sich zum Patienten wie ein Barcode
  zur Tafel Schokolade: 
  Niemand kauft eine Packung wegen des Codes.}
\end{Synopsis}

\begin{figure}[bh]
\begin{center}
\includegraphics[width=0.75\linewidth]{./images/wikipedia-pd/van_Gogh_Starry_Night_Over_the_Rhone.jpg}

\Captionof{figure}{Krank? \SourcePicture{Vincent van Gogh: Sternennacht über der
Rhone (September 1888).}}
\end{center}
\end{figure}


Der Krankheitswert eines Zustandes beginnt dort, wo dieser
Zustand für den Patienten selbst oder für seine Umgebung ein
Problem darstellt. Dabei gibt es keine scharfe Grenze und
die Grauzone ist beträchtlich.

\Text{Psychiatrische Erkrankung in der Praxis}
{%
In der Gesellschaft (repräsentiert durch die Gesetzgebung)
gibt es die Übereinkunft, dass bestimmte Krankheitsbilder
bzw. Symptome als behandlungswürdig bzw. sogar
behandlungspflichtig angesehen werden. Als
behandlungspflichtig gelten zum Beispiel die \E{Eigen-} und
\E{Fremdgefährdung} aufgrund einer psychiatrischen
Erkrankung, hier hat der Gesetzgeber Maßnahmen nach dem
\T{Unterbringungsgesetz} (\T{UbG}) vorgesehen. Auch der
geistige Verfall (z.\,B. im Rahmen einer (Alters-)Demenz)
kann rechtliche Schritte nach sich ziehen, im Sinne einer
Besachwaltung.
}{←}{}{}{}



\TextSlide{Der psychiatrische Patient ist nicht entrechtet!}
{%
Es ist hierbei jedoch festzuhalten, dass auch ein
psychiatrischer Patient -- genauso wie jeder andere Patient
auch -- über die selben Selbstbestimmungsrechte verfügt,
d.h. es ihm grundsätzlich zusteht eine Behandlung zu
verweigern. Sollten im Rahmen der gesetzlichen Bestimmungen
(\EE{Unterbringungsgesetz}) Maßnahmen durchzuführen sein, so
sind diese von der \EE{Exekutive} (Polizei) vorzunehmen. Das
Ziel des Unterbringungsgesetz ist es, das Maß der
\E{Einschränkung der Freiheit des Patienten so gering wie
möglich zu halten}.

Auch im Rahmen der
\EE{Besachwaltung} kann ein \E{Gericht} das Selbstbestimmungsrecht des
Patienten einschränken. 
}{
\begin{itemize}
  \item Der psychiatrische Patient hat grundsätzlich die gleichen Rechte!
  \item Besonderheiten:
  \begin{itemize}
    \item Unterbringungsgesetz (UbG)
    \item Besachwaltung
  \end{itemize}
\end{itemize}
}{}{}{}


\begin{Synopsis}
  \SynopsisItem{Der Rettungsdienst führt von sich aus keine
  Zwangsmaßnahmen bei psychiatrischen Patienten durch.}
\end{Synopsis}

Für die Arbeit im Rettungsdienst ist die Kenntnis einiger
spezieller Krankheitsbilder und Symptome notwendig, sowie
ein allgemeines Verständnis über den Umgang mit
psychiatrischen Patienten.

\section{\HeadingKeyword{Umgang} mit psychiatrischen Patienten}
\label{topic:umgang-psy}

Wie in der Einleitung erwähnt, ist eine psychiatrische
Erkrankung etwas durchaus Alltägliches, oft stellt selbige
nicht einmal den Behandlungs- bzw. Transportgrund dar.
Ebenso wurde erwähnt, dass der psychiatrische Patient
grundsätzlich über die selben (Selbstbestimmungs-)Rechte
wie jeder andere Patient verfügt. 

\paragraph{Haltung gegenüber dem Patienten} Dementsprechend
gelten grundsätzlich auch die gleichen Regeln für die
Gesprächsführung wie sie unter
\prettyref{topic:kommunikationsregeln} besprochen wurden. Der
Patient verdient den gleichen Respekt und die gleiche
Wertschätzung wie ein körperlich Erkrankter.

\paragraph{Auftreten} 

\begin{itemize}
    \item Besonders wichtig ist ein \EE{ruhiges,
    sachliches Auftreten}.
    \item  Eventuelle \EE{Aufregung} (durch den
    Patienten; Angehörige; Feuerwehr; Polizei;
    Spezialeinsatzgruppe mit Rammbock, Tränengas,
    Maschinengewehr und Angst; \ldots) \EE{soll nicht auf
    den Helfer abfärben}.
    \item Der Patient soll sich \EE{nicht beengt},
    bedroht und -- wenn möglich -- nicht bevormundet
    \EE{fühlen}.
    \item Es sollen \EE{keine falschen oder nicht
    einlösbare Versprechen} gegeben werden.
   \item \EE{Keine Lügen!}
\end{itemize}

\paragraph{Eigenschutz} Auch wenn vom durchschnittlichen
psychiatrischen Patienten keine höhere Gefahr als von der
Normalbevölkerung ausgeht, so ist es doch ratsam
grundlegende Punkte des Eigenschutzes zu beachten:

\begin{enumerate}
    \item Fluchtweg freihalten
    \item Rücken freihalten
    \item Patienten nicht beengen/bedrängen
    \item Patienten über Handlungen informieren und evtl.
    extra um Einverständnis fragen (angurten, \ldots)
\end{enumerate}


Diese Punkte sind grundsätzlich für jeden Einsatz sinnvoll!



\section{\HeadingKeyword{Symptome}}

\subsection{Der Patient mit \HeadingKeyword{Wahnvorstellungen}}

Menschen können -- aus unterschiedlichen Gründen --
Wahnvorstellungen (Halluzinationen) entwickeln. Diese
können sehr unterschiedlich sein und alle Sinne betreffen.
Man kann Dinge \E{hören} (Stimmen hören, \ldots)  oder
Sachen \E{sehen} (\enquote{weiße Mäuse}, \ldots) die nicht
existieren. Der Betroffene nimmt dies jedoch als vollkommen
\enquote{echt} und real wahr. 

Der Wirklichkeitsbezug kann teilweise erhalten oder nahezu
komplett verloren gegangen sein. 



\minisec{Fallgeschichte}

\begin{quotation}
Eine 37-jährige Patientin kommt aufgrund eines
grippalen Infektes während der Wintermonate in die
Ordination eines Arztes für Allgemeinmedizin. Während des
Gespräches erzählt die Patienten, dass sie schon seit
mehreren Jahren arbeitsunfähig sei, da sie an
Halluzinationen leide. Jedes Mal wenn sie das Badezimmer
betritt, sieht sie eine ihr unbekannte Frau beim Fenster
hinein schauen, obwohl sie im zweiten Stock wohnt. Sie
erzählt, dass sie wisse, dass es nicht real sein kann,
trotzdem nimmt sie diese Frau so wahr als würde sie
tatsächlich vor dem Fenster stehen. Sie ist nun schon seit
Jahren in psychiatrischer Behandlung, allerdings stellt
sich keine Besserung ihrer Symptome ein. Die Patientin gibt
an einen massiven Leidensdruck zu haben.
\end{quotation}



\paragraph{Fazit}

Auch psychiatrische Patienten haben \enquote{normale}
Alltagsprobleme, wie z.\,B. einen grippalen Infekt. Die
Wahnvorstellungen werden von den Patienten komplett real
erlebt, auch wenn sie vielleicht sogar wissen, dass sie nicht
real sein können. 

Diese Patientin hat noch eine teilweise
intakte Realitätswahrnehmung: Psychiatrische Erkrankungen
oder Symptome funktionieren \E{nicht} nach dem
\enquote{Alles-oder-Nichts}-Prinzip.


\paragraph{Wahnidee -- Wahnsysteme}

Auch die Wahrnehmung der Umgebung insgesamt kann gestört
sein, auch hier gibt es verschiedene Abstufungen. 

Die einfachste Form ist die \T{Wahnidee}. Hierbei hat der
Patient eine klar abgegrenzte Vorstellung von etwas, was mit
der Realität nicht vereinbar ist. 

\begin{quotation}
\emph{\enquote{Das
Gebühren-Info-Service des ORF kontrolliert jede Woche am
Verteilerschrank, ob das Antennenkabel korrekt vom
Verteiler abgesteckt wurde.}}
\end{quotation}

Diese Wahnideen können in weiterer Folge zu einem
\T{Wahnsystem} ausgebaut werden 



\begin{quotation}
\Speech{\enquote{Das Gebühren-Info-Service
des ORF ist eine Deckmantelorganisation des israelischen
Geheimdienstes Mossad, welcher im Auftrag der CIA geheime
Abhöranlagen im ganzen Wohnhaus installiert hat. Bei
unliebsamen Äußerungen der Bewohner werden Strahlen, welche
das Denken beeinflussen können, in die Wohnungen geleitet.
Um diesen zu entgehen hat der Patient das Antennenkabel am
Verteilerkasten abgesteckt und zahlt keine Fernsehgebühren
mehr. Da ihn nun der Geheimdienst nicht mehr erreichen
kann, haben sie vom Gebühren-Info-Service einen Agenten
vorbei geschickt, welcher nun gewaltsam die Gedanken des
Patienten wegnehmen wollte.}}
\end{quotation}

\paragraph{Wahnthemen} \Z{Ein Wahn kann unterschiedliche
Themen haben, die häufigsten sind z.\,B. Verfolgungswahn,
Größenwahn oder der Verarmungswahn.}

% M %%%%%%%%%%%%%%%%%%%%%%%%%%%%%%%%%%%%%%%%%%%%%%%%%%% M %
% Wahnvorstellungen
\ProcedurePrintDefault{PF22091C}

% \subsubsection{Realitätsverlust}
% 
% \paragraph{Mögliche Ursachen}
% 
% Endogene Psychosen (Schizophrenie,  schizoaffektive Psychose)
% 
% Organische Psychosen (Alkoholparanoia,  Drogenpsychose (z.\,B. Coke bugs
% ),  Encephalitis)
% 
% Realitätsverlust - Wahn
% 
% Wahn = falsche, allerdings in sich  logische Auffassung der Realität,
% an der  unkorrigierbar festgehalten wird  (paranoid = wahnhaft) 
% 
% Themen: Verfolgungswahn, Größenwahn,  Verarmungswahn,...


% \subsubsection{Realitätsverlust - Halluzination}
% 
% 
% Halluzinationen: Trugwahrnehmungen,  für alle Sinne möglich
% (weiße Mäuse,  Stimmen,
% \enquote{Geisterberührungen},...)
% 
% 
% Psychomotorische Unruhe /  Aggressivität
% 
% 
% Organische Ursachen ausschließen!!!
% 
% {\mdseries\color[rgb]{0.0,0.0,0.5019608}
% Intoxikation durch Alkohol, Drogen,  Medikamente}
% 
% 
% Hypoglykämie
% 
% St.p. SHT
% 
% Hirntumor
% 
% psychomotorische Unruhe /  Aggressivität
% 
% Schwere Erregungszustände:
% 
% Katatonie : Starre mit Sprechstereotypien, Grimassieren,
% Befehlsautomatismen, Patient imitiert Untersucher
% 
% Manie: Distanzlosigkeit, aggressiv gefärbte Agitiertheit
% 
% psychoreaktiv: abnorme Persönlichkeit
% 
% psychosoziale Krise: Reaktion auf Anlaß / Streßreaktion
% 
% psychomotorische Unruhe /  Aggressivität
% 
% Hysterie (funktionelle Erregungszustände):
% 
% anlaßgebunden, entsprechende Grundpersönlichkeit
% 
% psychosomatischer Funktionsverlust (Lähmungen, Erblindung) ohne echte
% organische Ursache
% 
% Entspricht psychischer Abwehrreaktion
% 
% \subsubsection{hysterischer Krampfanfall:}
% 
% {\mdseries
% keine Verletzungen, keine Apnoe, atypische Krämpfe, kein
% Stuhl-/Harnverlust}
% 
% {\mdseries
% unbewußt \enquote{gespielt}}



% \subsection{\Z{Reserviert}\A{Der \emph{stuporöse}
% Patient}}

%\Definition{Stupor: Erschlaffung}
% 
% \paragraph{\TxSymptome}
% 
% \begin{itemize}
%     \item ausdrucksloser Patient
%     \item unbeweglich
%     \item starre Mimik
%     \item spricht nicht, ißt nicht ={\textgreater} ev. bis Verhungern!
%     \item Suchterkrankungen
%     \item primär psychisches Problem
%     \item sekundär körperliche und soziale Folgen
%     \item eigengesetzlicher Ablauf
%     \item zunehmender Kontrollverlust über  eigenes Verhalten
% \end{itemize}

% \subsubsection{Manie}
% 
% {\mdseries
% rastlose Unruhe - \enquote{energiegeladen}}
% 
% {\mdseries
% gereiztes Verhalten}
% 
% {\mdseries
% Überaktivität}
% 
% {\mdseries
% Distanzverlust}
% 
% {\mdseries
% Verwirrungszustände}
% 
% {\mdseries
% Aggression}
% 
% {\mdseries
% Selbstüberschätzung}
% 
% {\mdseries
% ev. als bipolare Störung abwechselnd mit  Depression}


\subsection{Psychomotorische \HeadingKeyword{Unruhe} und \HeadingKeyword{Aggressivität}}

Unruhe und Aggressivität kann ein Anzeichen für viele
Krankheitsbilder sein. Die Differentialdiagnosen umfassen
z.\,B. Blutzuckerstörungen, Demenz, Intoxikation,
Panikattacken, Manie, akute Psychosen u.\,v.\,a.\,m. Hier soll nur
auf die grundsätzlich zu bedenkenden Maßnahmen im Umgang mit diesen Patienten
eingegangen werden. 

% M %%%%%%%%%%%%%%%%%%%%%%%%%%%%%%%%%%%%%%%%%%%%%%%%%%% M %
\ProcedurePrintDefault{PF92091C}

\subsection{\HeadingKeyword{Suizidalität}}
\label{topic:suizidalität}
\label{topic:suizid1}

\DefinitionInPara{\T{Suizid} bezeichnet die absichtliche
Selbsttötung. Als \T[Suizid!erweiterter]{erweiterter Suizid} wird die Tötung der eigenen und fremder Personen bezeichnet.}
Selbstmorde und Selbstmordversuche (\T{SMV}) gehören zu den
häufigen Einsatzgründen im Rettungsdienst. Häufig sind
Suizid- oder sonstige selbstschädliche Handlungen nicht auf
den ersten Blick als solche erkennbar. 

Wenn der Verdacht auf eine suizidale Handlung besteht, soll
dies vertraulich, aber \EE{direkt angesprochen und geklärt
werden}. Es ist dabei wichtig, einen \Ca{neutralen Tonfall}
zu verwenden und \Ca{nicht anklagend zu klingen!}

\begin{quotation}
\Speech{\enquote{Wollten Sie sich verletzen?}}

\Speech{\enquote{Wollten Sie sich umbringen?}}
\end{quotation}

Häufig berichten Patienten über Gefühle der
Hoffnungslosigkeit, Einsamkeit, Verzweiflung, innerer Leere
oder lassen einen fehlenden Realitätsbezug erkennen.
\EE{Die geäußerten Aussagen und Gefühle des Patienten sollen nicht
gewertet, verharmlost oder verniedlicht werden!}

Ein misslungener -- oft zögerlicher -- Selbstmordversuch
wird oft von Unwissenden im Sinne von \emph{\enquote{Der hat es gar nicht so gemeint}}
aufgefasst. Das ist falsch\footnote{Zumindest meistens.}!
Ein Selbstmordversuch ist oft der letzte Hilferuf, den ein
Mensch -- sei es aus Vereinsamung oder aus anderen Gründen
-- aussenden kann. 

\begin{Danger}
  \item Jede Suizidäußerung des Patienten muss ernst
  genommen und dokumentiert werden.
  \item Jeder Selbstmordversuch muss ernst genommen
  werden, auch wenn er noch so lächerlich oder unglaubwürdig
  aussieht.
\end{Danger}


\subsection{Der \HeadingKeyword{unmittelbar eigen- oder
fremdgefährdende} Patient}
\label{topic:unterbringung}
\label{topic:suizid2}


\Text{{\TxQuerverweise}}{%
\begin{itemize}
  \item Unterbringung: \dotfill \FullPrettyRef{topic:unterbringung-jus}
\end{itemize}
}{}{}{}{}





\bigskip

\section{\HeadingKeyword{Krankheitsbilder}}

\subsection{\HeadingKeyword{Psychose}}
\index{Psychose}
\label{topic:psychose}
\TODO{EMHO}{2010-08-04}{Psychose: Text fehlt}{}

\TextSlide{{\TxBeschreibung}: Psychose}
{\DefinitionInPara{\T{Psychose} ist eine allgemeine Bezeichnung für eine \E{psychische
Störung} mit strukturellem \E{Wandel des Erlebens}
\cite{Pschy259}.}
Man unterscheidet zwischen organischen, körperlich
begründbaren, und nicht-organischen, d.\,h. körperlich nicht
begründbaren Psychosen. In beiden Gruppen kann man weitere
Unterscheidungen treffen, vgl.
\prettyref{tab:psychosen-einteilung}. Körperliche Ursachen für Psychosen
können Störungen und Verletzungen des ZNS, z.\,B. bei
Hirntumoren, Folge eines SHT, Intoxikationen,
Stoffwechselstörungen u.\,v.\,a.\,m. sein.}
{%
\begin{itemize}
  \item Psychische Störung mit strukturellem Wandel des Erlebens
\end{itemize}
}


\begin{Synopsis}
  \SynopsisItem{Grundsätzlich muss bis zum Beweis des Gegenteils von
  einer \EE{organischen} Psychose (mit behebbarer Ursache)
  ausgegangen werden!}
\end{Synopsis}

% \TableWide
% 	{<Titel>}%1
% 	{<Beschreibung>}%2
% 	{<Label>}%3
% 	{<Quelle>}%4
% 	{<Lizenz>}%5
% 	{<Tabelle>}%6
% 	{<Float>}%8
\TableWide
{Einteilung der Psychosen nach \cite{Pschy259}}%1
{}%2
{\label{tab:psychosen-einteilung}}%3
{}%4
{}%5
{%
\Z{%
\begin{tabularx}{\linewidth}{XX}
\toprule\addlinespace
\multicolumn{2}{c}{\TabHeading{Psychosen}}
\\
\multicolumn{1}{c}{\TabHeading{Organische}} & \multicolumn{1}{c}{\TabHeading{Nicht-Organische}}
\\\addlinespace\cmidrule(lr){1-1}\cmidrule(lr){2-2}
\begin{enumerate}
  \item Akute organische Pschosen
  \begin{enumerate}
    \item Delir
    
    z.\,B. Alkoholentzugsdelir
    (\FullPrettyRef{topic:alkoholentzugsdelir})
    \item Dämmerzustand 
    \item Durchgangssysndrom 
    
    z.\,B. \FullPrettyRef{topic:durchgangssyndrom}
  \end{enumerate}
  \item Chronische organische Psychosen
  \begin{enumerate}
    \item frühkindliches organisches Psychosyndrom
    \item Hirndiffuses Psychosyndrom
    \item Hirnlokales Psychosyndrom
  \end{enumerate}
\end{enumerate}
&
\begin{enumerate}
  \item Schizophrene Psychosen
  \item Affektive Psychosen
  \item Schizoaffektive Psychosen
\end{enumerate}
\\\bottomrule
\end{tabularx}
}
}%6
{H}%8


\Text{\TxSymptome}
{Die Symptome einer Psychose sind vielfältig und abhängig
von der konkreten Störung und dem Schweregrad.
Sie sind durch \EE{Störungen  des Denkens, der Wahrnehmung
und des Affektes} charaktersisiert
\cite{ScholzNotfallmedizin2-20.1-PsychiatrischeNot}. 

}
{}
{}
{}
{}


% \clearpage
\subsection{\HeadingKeyword{Demenz}}



\TextSlide{{\TxBeschreibung}}{%
\DefinitionInPara{Die \T{Demenz} ist eine fortschreitende Einschränkung der  geistigen
Leistungsfähigkeit des Gehirns.} Die häufigste Form ist die
\I{Altersdemenz}, daneben gibt es noch viel andere Formen der Demenz. Die bekannteste Form ist die Demenz im Rahmen der \I{Alzheimer}-Krankheit. %
\index{Demenz!Alters-}\index{Demenz!Alzheimer-}%
Die geistige Leistungseinschränkung betrifft alle höheren Hirnfunktion, z.\,B. das Gedächtnis, das Denkvermögen, die Sprache, die Motorik, sowie die Persönlichkeitsstruktur und Affektverarbeitung.
}{
\begin{itemize}
  \item Geistige Leistungseinschränkung
  \item Viele Formen
  \item Alle höheren Hirnfunktionen können betroffen sein
\end{itemize}
}{}{}{}

\TextSlide{\TxSymptome}{%
Die Symptome sind je nach Form, 
Fortschritt 
und Person vielfältig 
und unterschiedlich, 
und darüber hinaus auch sehr wechselhaft 
und von der Tagesverfassung abhängig. 
%
Klassisch ist eine zeitliche oder örtliche Desorientiertheit,
Gedächtnisschwäche 
(insbesonders das Kurzzeitgedächtnis betreffend), 
das Verkennen von Personen, 
oder auch psychomotorische Unruhe und Aggressivität. 
}{
\begin{itemize}
    \item Zeitliche oder örtliche Desorientierung
    \item Gedächtnisschwäche
    \item Psychomotorische Unruhe, Aggressivität
    \item Personenverkennung
    \item Wechselnde Symptomatik!
\end{itemize}
}{}{}{}



\subsection{\HeadingKeyword{Depression} und \HeadingKeyword{Manie}}
\label{topic:manie}
\label{topic:depression}

\TextSlide{{\TxBeschreibung}: Depression und Manie}
{%
\DefinitionInPara{Die \T{Depression} ist eine sehr häufige psychiatrische
Erkrankung, welche durch gedrückte Stimmung,
Interessensverlust, Freudlosigkeit, Antriebsminderung,
erhöhte Ermüdbarkeit und Aktivitätseinschränkung
charakterisiert ist \cite{ICD10De}.
Die \T{Manie} ist dagegen ein Überbegriff für Auffälligkeiten des
Affektes, der Antriebs- und Willenssphäre sowie des
Denkens \cite{ScholzNotfallmedizin2-20.1-PsychiatrischeNot}.}

Die Depression ist eine sehr häufige Erkrankung
(\enquote*{Volkskrankheit}), etwa 10\PC* der Patienten einer
allgemeinmedizinischen Praxis leiden unter ihr
\cite{WunnPsych1}. Sie tritt mit einer Vielzahl von
unterschiedlichen Symptomen auf und ist oft nur schwer von
alltäglicher Verstimmung abzugrenzen
\cite{ScholzNotfallmedizin2-20.1-PsychiatrischeNot}.

Die Depression tritt oft episodenhaft auf. Bei den typischen
Episoden leidet der betroffene Patient unter einer
\EE{gedrückten Stimmung} und einer \EE{Verminderung von
Antrieb und Aktivität}. Die Fähigkeit zu Freude, das
Interesse und die Konzentration sind vermindert. Ausgeprägte
\EE{Müdigkeit} kann nach jeder kleinsten Anstrengung
auftreten. Der Schlaf ist meist gestört, der Appetit
vermindert. Selbstwertgefühl und Selbstvertrauen sind fast
immer beeinträchtigt. Sogar bei der leichten Form kommen
Schuldgefühle oder Gedanken über eigene \EE{Wertlosigkeit}
vor.

Die gedrückte Stimmung verändert sich von Tag zu Tag wenig,
\EE{reagiert nicht auf Lebensumstände} und kann von so
genannten \enquote{somatischen} Symptomen begleitet werden, wie
Interessenverlust oder Verlust der Freude, Früherwachen,
Morgentief, deutliche psychomotorische Hemmung,
Agitiertheit, Appetitverlust, Gewichtsverlust und
Libidoverlust. Abhängig von Anzahl und Schwere der Symptome
ist eine depressive Episode als leicht, mittelgradig oder
schwer zu bezeichnen \cite{ICD10De}. 
Im schweren Fällen kann es zum Übergang in ein \E{suizidales Syndrom} kommen (\FullPrettyRef{topic:suizidalität}).

Die Episoden können regelmäßig wiederkommen (rezidivierende
depressive Störung). Die erste Episode kann in jedem Alter
zwischen Kindheit und Senium auftreten, der Beginn kann akut
oder schleichend sein, die Dauer reicht von wenigen Wochen
bis zu vielen Monaten.
\cite{ICD10De}
}{%
\begin{itemize}
  \item Depression:
  \begin{itemize}
    \item Gedrückte Stimmung, Interessensverlust, Antriebslosigkeit
    \item Sehr häufig
    \item Unterschiedliche Schweregrade
  \end{itemize}
  \item Manie:
  \begin{itemize}
    \item Auffällige Antriebs- und Affektsteigerung
  \end{itemize}
  \item Suizidalität in schweren Fällen möglich
\end{itemize}
}







\subsection{\HeadingKeyword{Alkohol-} und 
\HeadingKeyword{Drogenentzug}}
\label{topic:alkohol-entzug}



\Text{\TxSymptome}
{%
\begin{itemize}
  \item Schwitzen
  \item Unruhe
  \item Zittern (Tremor)
  \item Delirium  (Bewusstseinsstörungen, Halluzinationen )
  \item Krampfanfälle
  \item RR-Entgleisung
  \item \E{\enquote{Cold turkey}} bei Heroinentzug
\end{itemize}
}{

}{}{}{}

\Text{Anamnese}
{
\begin{itemize}
  \item Plötzlich kein Zugang mehr zu Alkohol:
  \item Kürzlicher Einzug in Altersheim
  \item kürzlich pflegebedürftig geworden
  \item \ldots
\end{itemize}
}{}{}{}{}












%%%%%%%%%%%%%%%%%%%%%%%%%%%%%%%%%%%%%%%%%%%%%%%%%%%%%%%%%%%
\label{topic:alkoholentzugsdelir}
\label{topic:durchgangssyndrom}
\paragraph{Entzugsdelir und Durchgangssyndrom} Im Extremfall
leidet der Patient an einem Entzugsdelir (Delirium tremens). Dabei kann  es zu
Zittern, Unruhe, Angst, \EE{Halluzinationen} (\enquote{weiße
Mäuse}, \ldots), vegetativer Entgleisung (Tachykardie,
Hypo-/Hypertonie) und zerebralen Krampfanfällen kommen. Es
handelt sich um einen \Ca{lebensbedrohlichen Zustand}!



\subsection{\HeadingKeyword{Panikattacke}}

Eine Panikattacke kann sich sehr unterschiedlich äußern und
körperliche Beschwerden vortäuschen.

\begin{itemize}
    \item Plötzliches Angstgefühl
    \item Dauer: Sekunden bis Minuten
    \item Thoraxbeschwerden (Infarktangst)
    \item Hyperventilation \index{Hyperventilation!bei Panikattacke}
    \item Mehrfach durchuntersuchte  Patienten, wirken gesund
    \item Erwartungsangst
    \item Vermeidungsverhalten
    \item Antriebslosigkeit / Stupor
\end{itemize}




\begin{Literary}

\fullcite{ThunRedenI46}

\fullcite{ThunRedenII30}

\fullcite{ThunRedenIII19}

\fullcite{ArgelanderErstinterview8}

\fullcite{WunnPsych1}

\fullcite{BattegayHandwPsych2}

\fullcite{ScholzNotfallmedizin2}
\end{Literary}


\ChapterEnd